\documentclass[10pt,a4paper]{amsart}
\usepackage[margin=19mm]{geometry}
\usepackage{amsmath,amssymb,mathtools}
\usepackage[scr=rsfs,cal=euler]{mathalfa}
\usepackage{xfrac}
\usepackage[unicode,hidelinks,bookmarks=false]{hyperref}
\newcommand{\ie}{i.e.\ }
\newcommand{\cf}{cf.\ }
\newcommand{\resp}{respectively\ }
\newcommand{\Subsection}{\subsection}
\providecommand{\nobreakdash}{\nobreak\hbox{-}}
\setlength{\emergencystretch}{2em}
\multlinegap0pt
\usepackage{lmodern}
\newtheorem{theorem}{Theorem}[section]
\newtheorem{lemma}[theorem]{Lemma}
\newtheorem{proposition}[theorem]{Proposition}
\newtheorem{corollary}[theorem]{Corollary}
\newtheorem{definition}[theorem]{Definition}
\newtheorem{remark}[theorem]{Remark}
\newcommand{\C}{\mathcal C}
\newcommand{\Hom}{\operatorname{Hom}}
\newcommand{\Lie}{\operatorname{Lie}}
\newcommand{\Const}{\operatorname{Const}}
\newcommand{\LL}{\widehat{\mathbb L}}
\newcommand{\id}{\operatorname{id}}
\begin{document}
\title[Continuous Comparison of Free Simplicial Prounipotent and Pr]{Continuous Comparison of Free Simplicial Prounipotent and Pro-$p$ Resolutions}
\author{Andrey Mikhovich}
\date{}
\maketitle
\noindent\emph{Source and licence.} Continuous Comparison of Free Simplicial Prounipotent and Pro-$p$ Resolutions. Version arXiv:2609.21513v1 (CC BY 4.0). This material is distributed under the Creative Commons Attribution 4.0 International licence, http://creativecommons.org/licenses/by/4.0/, as recorded in the arXiv OAI licence metadata for this exact version and shown on the abstract page https://arxiv.org/abs/2609.21513v1. Full rights notice: licenses/arxiv-2609.21513-CC-BY-4.0.txt.\par\smallskip
\noindent\textit{Editorial scope.} Counted unit: the original \texttt{theorem} environment \texttt{thm:comparison} at source lines 790--800 together with its complete proof at lines 801--836; the delivered document keeps source lines 88--837 so that every referenced label and every citation resolves inside it. Native editable source is the paper's own concatenated TeX; the AMS wrapper is editorial formatting only. No publisher class, upstream script or unrelated artwork is redistributed.
\begin{equation}\label{eq:partition}
 F(S)\cong\coprod_{j=1}^t F(S_j).
\end{equation}
Indeed, a convergent assignment on $S$ is precisely a convergent
assignment on each of its finitely many parts. This statement concerns
coproducts in $\C$, not abstract free products of underlying groups.

\section{Continuous lifting}

We first establish the input needed in all simplicial degrees.

\begin{lemma}\label{lem:linear}
A continuous surjective linear map $q:V\to W$ between linearly compact
$k$-spaces admits a continuous linear section.
\end{lemma}
\begin{proof}
Continuous duality identifies linearly compact spaces contravariantly with
discrete vector spaces. In particular $V\cong(V^\vee)^*$, where
$V^\vee=\Hom_{\rm cts}(V,k)$ and $(-)^*$ is the full dual with its
linearly compact topology. The map $q^\vee:W^\vee\hookrightarrow V^\vee$
is injective. Extend a basis of its image to a basis of $V^\vee$, and
choose a linear retraction $r:V^\vee\to W^\vee$.
Then $s=r^*:W\to V$ is continuous and $qs=\id_W$.
\end{proof}

For a vector space $V$, write $\mathbb L(V)$ for the ordinary free Lie
algebra on $V$.  For any Lie algebra $\mathfrak g$ its lower central series is
\[
 \gamma_1\mathfrak g=\mathfrak g,\qquad
 \gamma_{n+1}\mathfrak g=[\mathfrak g,\gamma_n\mathfrak g]
 \quad(n\geq1).
\]
Thus $\mathbb L(V)/\gamma_{c+1}\mathbb L(V)$ is the free nilpotent Lie
algebra of class at most $c$ on $V$.

For a linearly compact space $E$ define
\begin{equation}\label{eq:free-lie}
 \LL(E)=\varprojlim_{W,c}
 \mathbb L(E/W)/\gamma_{c+1}\mathbb L(E/W),
\end{equation}
where $W$ runs through open subspaces of finite codimension and $c\geq1$.
The inverse system is ordered by refinement: $(W',c')$ maps to $(W,c)$
whenever $W'\subseteq W$ and $c'\geq c$, via the evident quotient map.
Every displayed factor is a finite-dimensional nilpotent Lie algebra.
The canonical map $E\to\LL(E)$ is continuous, and its image
topologically generates $\LL(E)$.

\begin{lemma}\label{lem:universal}
For every pronilpotent Lie algebra $\mathfrak u$, restriction to $E$ gives
a bijection
\[
 \Hom_{\rm cts,Lie}(\LL(E),\mathfrak u)
 \cong \Hom_{\rm cts,lin}(E,\mathfrak u).
\]
\end{lemma}
\begin{proof}
Let $a:E\to\mathfrak u$ be continuous and linear.  Since
$\mathfrak u$ is pronilpotent, we may write
\[
 \mathfrak u\cong\varprojlim_J \mathfrak u/J,
\]
where $J$ runs through the open ideals for which $\mathfrak u/J$ is
finite-dimensional and nilpotent.  Fix such a $J$, and write
$a_J:E\to\mathfrak u/J$ for the composite of $a$ with the quotient map.
Because $\mathfrak u/J$ is finite-dimensional and discrete, continuity of
$a_J$ implies that $W:=\ker(a_J)$ is open; moreover $E/W$ is
finite-dimensional.  More generally one may replace $W$ by any open
finite-codimensional subspace contained in $\ker(a_J)$.  Thus $a_J$
factors as a linear map
\[
 \bar a_{J,W}:E/W\longrightarrow\mathfrak u/J.
\]
Choose $c$ at least the nilpotency class of $\mathfrak u/J$.  By the
ordinary universal property of the free Lie algebra, $\bar a_{J,W}$
extends uniquely to a Lie homomorphism $\mathbb L(E/W)\to\mathfrak u/J$.
Since every $(c+1)$-fold iterated bracket vanishes in $\mathfrak u/J$,
this homomorphism kills $\gamma_{c+1}\mathbb L(E/W)$ and therefore factors
uniquely through
\[
 \phi_{J,W,c}:\mathbb L(E/W)/\gamma_{c+1}\mathbb L(E/W)
       \longrightarrow\mathfrak u/J.
\]
Composing with the canonical projection
$p_{W,c}:\LL(E)\to
\mathbb L(E/W)/\gamma_{c+1}\mathbb L(E/W)$ gives a continuous Lie map
$\Phi_J:=\phi_{J,W,c}p_{W,c}:\LL(E)\to\mathfrak u/J$.

We now verify that $\Phi_J$ is independent of the auxiliary choices.
Suppose that $(W_1,c_1)$ and $(W_2,c_2)$ are two admissible choices for the
same $J$.  Put
\[
 W_3=W_1\cap W_2,\qquad c_3=\max\{c_1,c_2\}.
\]
Then $(W_3,c_3)$ refines both pairs.  After passage to the common factor
$\mathbb L(E/W_3)/\gamma_{c_3+1}$, both maps to $\mathfrak u/J$ extend
the same linear map $E/W_3\to\mathfrak u/J$ induced by $a_J$.
The free nilpotent universal property therefore makes them equal.
Consequently the two composites from $\LL(E)$ coincide.  This proves
independence of both $W$ and $c$.

Next let $J'\subseteq J$.  Denote by
$\rho_{J'J}:\mathfrak u/J'\to\mathfrak u/J$ the quotient map.  Choose a
common refinement $W$ of the subspaces used for $J'$ and $J$, and choose
$c$ at least the nilpotency classes of both quotients.  On the finite free
nilpotent factor $\mathbb L(E/W)/\gamma_{c+1}$, the two maps
\[
 \rho_{J'J}\Phi_{J'}
 \quad\text{and}\quad
 \Phi_J
\]
are induced by the same linear map $a_J:E\to\mathfrak u/J$; hence, again
by uniqueness of the free nilpotent extension, they are equal.  Thus the
family $(\Phi_J)_J$ is compatible with all transition maps as $J$ varies.

The universal property of the inverse limit now gives a unique Lie map
\[
 \Phi:\LL(E)\longrightarrow\varprojlim_J\mathfrak u/J
       \cong\mathfrak u
\]
whose $J$-component is $\Phi_J$.  It is continuous because the topology on
$\mathfrak u$ is the inverse-limit topology and each $\Phi_J$ is
continuous.  By construction $\Phi|_E=a$, so restriction is surjective.

Finally, if $\Psi:\LL(E)\to\mathfrak u$ is another continuous Lie map
with $\Psi|_E=a$, then for every $J$ the maps to $\mathfrak u/J$ obtained
from $\Psi$ and $\Phi$ agree on $E$.  They therefore agree on the Lie
subalgebra generated by $E$, which is dense in $\LL(E)$; continuity and
the Hausdorffness of the finite-dimensional quotient $\mathfrak u/J$ imply
that they agree on all of $\LL(E)$.  Since the quotient maps
$\mathfrak u\to\mathfrak u/J$ separate points, $\Psi=\Phi$.
This proves injectivity and completes the proof.
\end{proof}

\begin{proposition}[Projective lifting]\label{prop:projective}
Let $q:U\twoheadrightarrow V$ be an epimorphism in $\C$.
Every morphism $f:F(S)\to V$ admits a morphism
$\widetilde f:F(S)\to U$ with $q\widetilde f=f$.
There is no restriction on the cardinality of $S$.
\end{proposition}
\begin{proof}
In characteristic zero the assertion holds more generally with
$F(S)$ replaced by $\exp\LL(E)$ for any linearly compact $E$.
Let $q_*:\mathfrak u\twoheadrightarrow\mathfrak v$ and
$f_*:\LL(E)\to\mathfrak v$ be the corresponding Lie morphisms.
Choose the continuous linear section $s$ of $q_*$ furnished by
Lemma~\ref{lem:linear}. The continuous linear map
$s\circ f_*|_E$ extends by Lemma~\ref{lem:universal} to a continuous
Lie morphism $\widetilde f_*:\LL(E)\to\mathfrak u$.
The maps $q_*\widetilde f_*$ and $f_*$ agree on $E$, and hence agree
everywhere. Exponentiation proves the assertion. The section $s$ is
not required to preserve Lie brackets.

In the pro-$p$ category this is the classical projectivity theorem for
free pro-$p$ groups, valid in arbitrary rank
\cite[Section~7.7]{RibesZalesskii}. We recall how its finite lifting
criterion applies here. For a diagram $F_p(S)\to B\leftarrow A$ with
$A\twoheadrightarrow B$ a surjection of finite $p$-groups, the images
of all but finitely many generators in $B$ are $1$. Lift the finitely
many nonidentity images and send every remaining generator to $1$.
The resulting assignment is convergent and extends to a continuous
solution. The standard projectivity criterion for pro-$p$ groups says
that solving all such finite embedding problems is equivalent to lifting
along arbitrary epimorphisms of pro-$p$ groups
\cite[Section~7.6]{RibesZalesskii}. This yields the stated lifting.
\end{proof}

\begin{corollary}[Relative lifting]\label{cor:relative}
Suppose $B=A\amalg F(S)$ in $\C$, $q:U\twoheadrightarrow V$ is an
epimorphism, and $a:A\to U$, $f:B\to V$ satisfy $qa=f|_A$.
Then $f$ has a lift $\widetilde f:B\to U$ with $\widetilde f|_A=a$.
\end{corollary}
\begin{proof}
Lift $f|_{F(S)}$ by Proposition~\ref{prop:projective} and combine this
lift with $a$ using the coproduct universal property.
\end{proof}

\section{Free simplicial objects and relative extension}

An augmentation $X_\bullet\to G$ means an augmentation to
$\Const(G)$. It is useful to spell out the matching objects, since they
are the objects against which the inductive lifting is performed. Put
$Y_{-1}=H$ and regard an augmented simplicial object as a functor on the
augmented simplex category. For $n\geq0$ define the \emph{full matching
object}
\[
 M_nY=\varprojlim_{\substack{\delta:[m]\hookrightarrow[n]\\-1\leq m<n}}Y_m,
\]
where the limit is taken in $\C$ over all proper injective maps into
$[n]$. Thus $M_nY$ is the object of all proper faces of a formal
$n$--simplex which agree on their common subfaces. The natural matching
map
\[
 b_n:Y_n\longrightarrow M_nY
\]
is induced by the face maps (and, in degree zero, by the augmentation).
In particular
\[
 M_0Y=H,\qquad M_1Y=Y_0\times_HY_0,
\]
and, for $n\geq1$, the codimension-one faces determine the whole
matching object, giving the concrete description
\[
 M_nY=\{(y_0,\ldots,y_n)\in Y_{n-1}^{n+1}:
              d_i y_j=d_{j-1}y_i\text{ for }i<j\}.
\]
For $n=1$ the single compatibility equation means that the two elements
of $Y_0$ have the same image in $H$. For example,
\[
 M_2Y=\{(y_0,y_1,y_2)\in Y_1^3:
 d_0y_1=d_0y_0,\ d_0y_2=d_1y_0,\ d_1y_2=d_1y_1\}.
\]
Under the tuple description, $b_n=(d_0,\ldots,d_n)$. Thus surjectivity
of $b_n$ says precisely that every compatible \emph{full boundary}
admits an $n$--simplex filling it; this is stronger data than a single
horn-filling statement.

For compatibility with Keune's notation, put
\[
 Z_{-1}Y=H,\qquad Z_nY=M_{n+1}Y\quad(n\geq0),
\]
and write
\[
 d^{(-1)}=b_0:Y_0\longrightarrow H,\qquad
 d^{(n)}=b_{n+1}:Y_{n+1}\longrightarrow Z_nY\quad(n\geq0).
\]
Thus $Z_nY$ is the \emph{simplicial kernel}: it consists of compatible
full boundaries in degree $n+1$. We call $Y_\bullet\to H$ a
\emph{resolution} when every $d^{(n)}$, $n\geq-1$, is an epimorphism.

\begin{proposition}[Keune: simplicial kernels and homotopy]
\label{prop:keune-simplicial-kernel}
Let $Y_\bullet\to H$ be an augmented simplicial group. For every
$n\geq0$ there is a natural bijection of pointed sets
\[
 \widetilde\gamma_n:
 \pi_n(Y_\bullet\to H)\xrightarrow{\ \cong\ }
 Z_nY/\operatorname{im}d^{(n)}.
 \tag{SK}
\]
To specify both sides and their group structures, put
\[
 \begin{aligned}
 N_mY&=\bigcap_{i=0}^{m-1}\ker d_i &&(m\geq1),\\
 C_nY&=\bigcap_{i=0}^{n}\ker d_i &&(n\geq1),\\
 C_0Y&=\ker(Y_0\to H).
 \end{aligned}
\]
and write
\[
 B_nY=d_{n+1}(N_{n+1}Y),\qquad
 D_nY=\operatorname{im}d^{(n)}.
\]
Thus $\pi_n(Y_\bullet\to H)=C_nY/B_nY$ is the homotopy group
of the augmented Moore complex. The map in \textup{(SK)} is induced by
\[
 \gamma_n(c)=(1,\ldots,1,c)\in Z_nY,\qquad c\in C_nY.
\]
The quotient on the right of \textup{(SK)} means the set of left cosets
$zD_nY$, with distinguished element $D_nY$. It is equipped with the
group structure transported from $C_nY/B_nY$ along
$\widetilde\gamma_n$. Explicitly, every coset has a representative
$\gamma_n(c)$, and the multiplication is
\[
 \bigl(\gamma_n(c)D_nY\bigr)\star
 \bigl(\gamma_n(c')D_nY\bigr)
 =\gamma_n(cc')D_nY.
\]
With this multiplication, \textup{(SK)} is a natural isomorphism of
groups. No normality of $D_nY$ in $Z_nY$ is assumed, and the coset
projection from $Z_nY$ is not asserted to be a group homomorphism.
In the pro-$p$ and prounipotent cases, images are closed; the topology,
respectively the prounipotent group structure, on the right of
\textup{(SK)} is transported from the Moore quotient.
Consequently, $d^{(n)}$ is surjective exactly when the degree-$n$
augmented homotopy obstruction vanishes.
\end{proposition}

\begin{proof}
For the Moore-complex description of homotopy groups, see
Inassaridze \cite[Chapter~2, pp.~35--36, Definition~2.2]{Inassaridze}.
The full-boundary bijection is given in
\cite[Chapter~2, Proposition~2.14, pp.~41--42]{Inassaridze}, more
generally for augmented pseudosimplicial groups. That proposition
explicitly describes the right-hand side as a quotient set.
We give the passage to full-boundary classes explicitly, using Keune's
simplicial-kernel notation \cite{Keune}, to specify the meaning of the
quotient in \textup{(SK)}.

The simplicial identities give $B_nY\subseteq C_nY$.
Moreover, $N_{n+1}Y$ is normal in $Y_{n+1}$, and $d_{n+1}$ is
surjective because $d_{n+1}s_n=\id$. Hence $B_nY$ is normal in
$Y_n$, and in particular in $C_nY$.
We claim that
\begin{equation}\label{eq:coset-comparison}
 \begin{aligned}
 \gamma_n(C_nY)\cap D_nY&=\gamma_n(B_nY),\\
 Z_nY&=\gamma_n(C_nY)D_nY.
 \end{aligned}
\end{equation}
Indeed, $\gamma_n(c)=d^{(n)}x$ means precisely that
$d_0x=\cdots=d_nx=1$ and $d_{n+1}x=c$. This is equivalent to
$x\in N_{n+1}Y$ with Moore boundary $c$, proving the first equality.

For the second equality, let $z=(z_0,\ldots,z_{n+1})\in Z_nY$.
Fill the horn consisting of its first $n+1$ faces: there is an
$x\in Y_{n+1}$ with $d_ix=z_i$ for $0\leq i\leq n$.
For completeness, such an $x$ is obtained by the finite recursion
\[
 \begin{aligned}
 x^{(0)}&=1,\qquad x=x^{(n+1)},\\
 x^{(i+1)}&=s_i\bigl(z_i(d_ix^{(i)})^{-1}\bigr)x^{(i)}
 \quad(0\leq i\leq n).
 \end{aligned}
\]
The matching equations ensure that each step preserves the earlier
faces. Put $c=z_{n+1}(d_{n+1}x)^{-1}$. The same equations imply
$c\in C_nY$, including the augmentation condition when $n=0$.
Thus
\[
 z\bigl(d^{(n)}x\bigr)^{-1}=\gamma_n(c),
\]
which proves the second equality in \eqref{eq:coset-comparison}.

It follows that $cB_nY\mapsto\gamma_n(c)D_nY$ is a well-defined
bijection. Since $B_nY$ is normal in $C_nY$, the displayed
multiplication $\star$ is independent of representatives and makes
this bijection a group isomorphism. The construction commutes with
morphisms of augmented simplicial groups, so the isomorphism is
natural. In particular the coset set is a singleton exactly when
$D_nY=Z_nY$, giving the asserted surjectivity criterion.

All operations in the horn-filling recursion are continuous in the
pro-$p$ case and are morphisms of underlying schemes in the
prounipotent case. Images are closed by compactness for pro-$p$
groups, and by closedness of continuous linear images of linearly
compact spaces after passing to Lie algebras in the prounipotent
case. Transporting the structure from the Moore quotient gives the
stated interpretation in both categories. For pro-$p$ groups the
transported topology also agrees with the coset topology, since the
induced continuous bijection is between compact Hausdorff spaces.
\end{proof}

Thus the resolving condition is not merely terminology: it says that
every class in the simplicial kernel is an actual full boundary. By
Proposition~\ref{prop:keune-simplicial-kernel}, for a general augmented
simplicial group the obstruction to this assertion is precisely a
homotopy class.

For $0\leq m\leq n$ write
\[
 D(n,m)=\operatorname{Sur}_{\Delta}([n],[m])
\]
for the set of actual order-preserving surjections in the simplex
category. This notation is important: an element of $D(n,m)$ is a
morphism of $\Delta$, not a formal word in elementary degeneracy
operators. Hence different degeneracy words which are identified by the
simplicial identities represent one and the same element of $D(n,m)$.
For example, the two words in
\[
 s_i s_j=s_{j+1}s_i\qquad(i\leq j)
\]
do not label two generators. Equivalently, one may start with formal
strings of degeneracies and quotient by the degeneracy identities; the
resulting equivalence classes are canonically parametrized by the
surjections in $D(n,m)$.

\begin{definition}\label{def:free}
A \emph{based free simplicial object} is a simplicial object together
with sets $S_m$ of nondegenerate free generators and specified
isomorphisms
\begin{equation}\label{eq:basis}
 X_n\cong F(B_n),\qquad
 B_n=\coprod_{0\leq m\leq n}\ \coprod_{\sigma\in D(n,m)}S_m^{\sigma}.
\end{equation}
Here $S_m^{\sigma}$ is one copy of $S_m$ for the actual surjection
$\sigma$; no additional copy is introduced for another word of
degeneracies representing the same $\sigma$. If
$\alpha_i:[n+1]\twoheadrightarrow[n]$ is the elementary surjection
corresponding to $s_i$, then
\[
 s_i(e,\sigma)=(e,\sigma\alpha_i).
\]
Thus the simplicial degeneracy identities hold on the basis literally,
because the corresponding composites are equal morphisms in $\Delta$.
All face and degeneracy maps are morphisms in $\C$. The object is of
\emph{finite type} if each $S_m$ is finite.
\end{definition}

This is the usual unique-degeneracy, or Eilenberg--Zilber, convention:
every basis element of $X_n$ is either a new element of $S_n^{\id}$ or
the unique degeneracy, indexed by an actual surjection $\sigma$, of a
nondegenerate generator in a lower degree. Thus ``free simplicial'' is
stronger than requiring each $X_n$ merely to be free. In each degree
\eqref{eq:partition} makes the new and degenerate families free factors.

\paragraph{Why a relative extension is used.}
Here the adjective \emph{relative} does not mean a group extension in the
short-exact-sequence sense.  It means that a simplicial subobject
$A_\bullet$ has already been constructed and is to be kept fixed while one
passes to a larger simplicial object $B_\bullet$ by freely adjoining new
nondegenerate generators, degree by degree, with prescribed compatible
full boundaries.  Thus the construction is ``relative to $A_\bullet$''.

This relative formulation is essential for the comparison argument rather
than merely convenient terminology.  In Lemma~\ref{lem:extension} a map
$a:A_\bullet\to Y_\bullet$ has already been prescribed and the required
lift on $B_\bullet$ must extend \emph{that same map}, not replace it by a
new choice.  Projectivity lifts the newly adjoined free factor, while the
relative coproduct decomposition combines that lift with the fixed map on
$A_\bullet$.  The same point is indispensable for homotopy uniqueness:
for the cylinder one starts with the two already prescribed endpoint maps
$f_0,f_1$ on $X_\bullet\amalg X_\bullet$ and extends them across
$C(X)_\bullet$.  Absolute degreewise freeness alone would not encode the
requirement that these previously constructed maps remain unchanged.  In
the topological categories considered here the relative formulation also
packages the infinitely many new generators into a single morphism in
$\C$, which is what preserves continuity.

We use exactly the same convention relatively. A \emph{relative free
extension} $A_\bullet\to B_\bullet$ is equipped with sets $T_m$ of new
nondegenerate relative generators and decompositions
\[
 B_n\cong A_n\amalg
 F\left(\coprod_{m\leq n}\ \coprod_{\sigma\in D(n,m)}T_m^\sigma\right).
\]
Again the index is the actual surjection $\sigma$, so generators related
by simplicial degeneracy identities are already identified and are not
adjoined independently. The faces of a new generator $t\in T_n$ lie in
the previously constructed degree $B_{n-1}$ and must be compatible:
\[
 (d_0t,\ldots,d_nt)\in M_nB.
\]
Equivalently, the attaching data in degree $n$ form a morphism from the
new free factor to the appropriate matching object of the lower
skeleton. Thus adjoining $T_n$ means adjoining genuinely new
nondegenerate generators together with one compatible full boundary for
each of them; all their degeneracies are then forced by the simplicial
identities.

\begin{lemma}[Simplicial extension]\label{lem:extension}
Let $A_\bullet\to B_\bullet$ be a relative free extension, with an
augmentation $B_\bullet\to G$. Let $Y_\bullet\to H$ be a resolution
and $\alpha:G\to H$ a morphism. Every simplicial morphism
$a:A_\bullet\to Y_\bullet$ over $\alpha$ extends to a simplicial
morphism $B_\bullet\to Y_\bullet$ over $\alpha$.
\end{lemma}
\begin{proof}
We give the induction in full.  Regard the augmentations as degree
$-1$ and put
\[
 u_{-1}=\alpha:G\longrightarrow H.
\]
We shall construct, for every $n\geq0$, maps
\[
 u_j:B_j\longrightarrow Y_j\qquad(-1\leq j\leq n)
\]
which form a morphism of augmented simplicial objects through degree
$n$, agree with $a_j$ on $A_j$, and are compatible with all face and
degeneracy operators whose source and target have degree at most $n$.

\emph{Base of the induction: $n=0$.}
In degree zero the relative free decomposition is
\[
 B_0\cong A_0\amalg F(T_0).
\]
Let $\varepsilon_B:B_0\to G$ be the augmentation and set
\[
 c_0=\alpha\varepsilon_B:B_0\longrightarrow H=M_0Y.
\]
Since the given map $a:A_\bullet\to Y_\bullet$ is over $\alpha$,
\[
 b_0a_0=c_0|_{A_0},
\]
where $b_0:Y_0\twoheadrightarrow H$ is the degree-zero matching map.
Apply Corollary~\ref{cor:relative} to
\[
 A_0\longrightarrow B_0=A_0\amalg F(T_0),
 \qquad b_0:Y_0\twoheadrightarrow H,
\]
with prescribed lift $a_0$ on $A_0$.  We obtain a morphism
\[
 u_0:B_0\longrightarrow Y_0
\]
extending $a_0$ and satisfying
\[
 b_0u_0=\alpha\varepsilon_B.
\]
Thus $u_{-1}=\alpha$ and $u_0$ form a morphism of augmented objects in
degrees $-1$ and $0$.  This is the induction base; no choice of
individual generators is being made.

\emph{Induction step.}
Let $n\geq1$ and suppose that
$u_{-1},u_0,\ldots,u_{n-1}$ have already been constructed and form a
morphism of the $(n-1)$--truncations.  Put
\[
 D_n=
 \coprod_{m<n}\ \coprod_{\sigma\in D(n,m)}T_m^\sigma,
 \qquad
 P_n=A_n\amalg F(D_n).
\]
Since the only element of $D(n,n)$ is the identity, the relative free
decomposition in degree $n$ can be rewritten as
\[
 B_n\cong P_n\amalg F(T_n).                 \tag{*}
\]
The factor $P_n$ is the part whose image is already forced by the lower
truncation.  Define
\[
 v_n:P_n\longrightarrow Y_n
\]
by $v_n|_{A_n}=a_n$ and, for $t\in T_m$, $m<n$, by
\[
 v_n(t^\sigma)=Y(\sigma)u_m(t),
 \qquad \sigma\in D(n,m).                  \tag{**}
\]
This is well defined because $\sigma$ is an actual surjection in
$\Delta$, rather than a chosen word in elementary degeneracies.  For a
fixed $n$ there are only finitely many possible surjections
$[n]\twoheadrightarrow[m]$.  The assignments in (**) are continuous on
each generator family, and the coproduct universal property together
with \eqref{eq:partition} therefore gives a morphism $v_n$ in $\C$.

We next verify explicitly that $v_n$ has exactly the boundary prescribed
by the already constructed lower-dimensional map.  It is enough to
check the codimension-one faces.  On $A_n$ this is true because $a$ is
simplicial.  Let $t\in T_m$, $m<n$, and let
$\sigma:[n]\twoheadrightarrow[m]$.  For the injection
$\delta_i:[n-1]\hookrightarrow[n]$ factor the monotone map
\[
 \sigma\delta_i:[n-1]\longrightarrow[m]
\]
uniquely as an epimorphism followed by a monomorphism,
\[
 \sigma\delta_i=\iota\rho,
 \qquad
 \rho:[n-1]\twoheadrightarrow[q],\quad
 \iota:[q]\hookrightarrow[m].             \tag{***}
\]
By the simplicial structure of the relative free extension,
\[
 d_i(t^\sigma)=B(\rho)B(\iota)t.
\]
Here $q\leq m<n$, so the induction hypothesis applies to every term in
this formula.  Consequently
\begin{align*}
 u_{n-1}d_i(t^\sigma)
 &=u_{n-1}B(\rho)B(\iota)t\\
 &=Y(\rho)u_qB(\iota)t\\
 &=Y(\rho)Y(\iota)u_m(t)\\
 &=Y(\sigma\delta_i)u_m(t)\\
 &=d_iY(\sigma)u_m(t)
  =d_iv_n(t^\sigma).
\end{align*}
If $q=m$, then $\iota=\id$ and the third equality is tautological; if
$q<m$, it is precisely the face compatibility already contained in the
induction hypothesis.  Thus all faces of $v_n$ agree with the maps
$u_j$ already constructed.  The same formula with an elementary
surjection in place of $\delta_i$ shows directly that
\[
 v_ns_i=s_iu_{n-1}
\]
on degree $n-1$.  Hence all degeneracy identities entering degree $n$
are also already satisfied on $P_n$.

Let
\[
 b_n^B:B_n\longrightarrow M_nB
\]
be the full matching map of the augmented simplicial object $B$.  Since
the maps $u_{-1},\ldots,u_{n-1}$ form a morphism of truncated augmented
simplicial objects, they induce a morphism of matching objects
\[
 M_n(u):M_nB\longrightarrow M_nY.
\]
Set
\[
 c_n=M_n(u)b_n^B:B_n\longrightarrow M_nY.
\]
The preceding boundary calculation says exactly that
\[
 b_n^Yv_n=c_n|_{P_n}.                       \tag{****}
\]
For $n=1$, condition (****) includes the equality of the two
augmentations in $H$; for $n>1$ the usual matching equations
$d_id_j=d_{j-1}d_i$ are already part of the induction hypothesis.

At this point the resolving hypothesis on $Y$ enters essentially.
For a new nondegenerate generator $t\in T_n$, the element
\[
 c_n(t)\in M_nY=Z_{n-1}Y
\]
is its already prescribed compatible full boundary. If $Y_\bullet\to H$
were merely an augmented simplicial group, there would be no reason for
$c_n(t)$ to have a filler in $Y_n$. In fact,
Proposition~\ref{prop:keune-simplicial-kernel} identifies the obstruction
to filling this boundary with the class
\[
 [c_n(t)]\in
 Z_{n-1}Y/\operatorname{im}d^{(n-1)}
 \cong \pi_{n-1}(Y_\bullet\to H)
 \qquad(n\geq1).
\]
Thus the simplicial kernel is exactly where the homotopy obstruction to
the inductive extension lives. For $n=0$ the analogous obstruction is
removed by the epimorphism $b_0:Y_0\twoheadrightarrow H$.

Because $Y_\bullet\to H$ is a resolution, all these obstructions
vanish simultaneously:
\[
 b_n^Y=d^{(n-1)}:Y_n\twoheadrightarrow
 M_nY=Z_{n-1}Y
\]
is an epimorphism.  We may therefore apply
Corollary~\ref{cor:relative} to the free factorization (*) with the map
$c_n$ and the prescribed lift $v_n$ on $P_n$.  Notice that we do not
choose unrelated fillers generator by generator: the relative
projective lifting produces one morphism in $\C$, so the required
continuity is preserved even for an infinite convergent family of new
generators. It gives a morphism
\[
 u_n:B_n\longrightarrow Y_n
\]
which extends $v_n$ and satisfies
\[
 b_n^Yu_n=c_n.                               \tag{*****}
\]
Equation (*****) gives all face identities in degree $n$.  Because
$u_n$ extends $v_n$, formula (**) gives all degeneracy identities from
degrees $<n$ into degree $n$.  Also $u_n|_{A_n}=a_n$.  Therefore
$u_{-1},u_0,\ldots,u_n$ form a morphism of augmented simplicial objects
through degree $n$.

Induction over $n$ produces a simplicial morphism
\[
 u:B_\bullet\longrightarrow Y_\bullet
\]
over $\alpha$ extending $a$.  Every lifting used above is a morphism in
$\C$ by Corollary~\ref{cor:relative}; hence continuity is preserved in
all degrees, including the case of infinitely many relative generators.

We emphasize the logical division of the proof. Relative freeness of
$A_\bullet\to B_\bullet$ and projectivity of the free factors provide
continuous lifts once a full boundary is known to be fillable. The
fact that \emph{every} compatible full boundary is fillable is exactly
the resolution hypothesis on $Y_\bullet\to H$; by
Proposition~\ref{prop:keune-simplicial-kernel}, without that hypothesis
the induction can stop on a nonzero homotopy obstruction.
\end{proof}

\section{The cylinder and the comparison theorem}

For a simplicial object $X$ define its cylinder by
\begin{equation}\label{eq:cylinder}
 C(X)_n=\coprod_{\tau:[n]\to[1]}X_n,
\end{equation}
where $\tau$ is order preserving. There are $n+2$ summands.
For a simplicial operator $\theta:[m]\to[n]$, the map from the
$\tau$-summand is $X(\theta)$ followed by the inclusion of the
$\tau\theta$-summand. This defines a simplicial object.
The two constant maps $[n]\to[1]$ give endpoint inclusions
$i_0,i_1:X\to C(X)$, and folding the summands gives $C(X)\to X$.
If $X$ is augmented to $G$, so is $C(X)$.

\begin{lemma}\label{lem:cylinder}
If $X$ is based free, the endpoint inclusion
$X\amalg X\to C(X)$ is a relative free extension.
\end{lemma}
\begin{proof}
By \eqref{eq:basis} and \eqref{eq:cylinder}, the free generators of
$C(X)_n$ have labels
\[
 (\sigma,e,\tau),\quad
 \sigma\in D(n,m),\quad e\in S_m,\quad
 \tau:[n]\to[1].
\]
Here again $\sigma$ is an actual surjection, so no duplicate labels arise
from different degeneracy words. A label is degenerate precisely when
the two monotone maps $\sigma$ and $\tau$ have a common repetition: there is an $i$ with
$\sigma(i)=\sigma(i+1)$ and $\tau(i)=\tau(i+1)$.
Collapse all maximal consecutive blocks on which the pair
$(\sigma,\tau)$ is constant. This gives a unique factorization
\[
 (\sigma,\tau)=(\bar\sigma,\bar\tau)\rho,
 \qquad \rho:[n]\twoheadrightarrow[r],
\]
where the reduced pair has no common repetition. The reduced label
is the unique nondegenerate label from which the original arises.
If $\tau$ is constant, its generators belong to the endpoint object,
and the reduced labels are exactly its ordinary nondegenerate ones.
All other reduced labels give the new relative generators.

For fixed $n$ there are only finitely many pairs of monotone maps
involved. Thus the partition into endpoints, new generators and
their degeneracies is a finite partition into copies of the $S_m$.
Each family retains its convergent-generator topology, and
\eqref{eq:partition} gives the required decomposition in $\C$.
Deleting a vertex sends a new generator to degree $n-1$, where its
image is already part of the lower skeleton. This proves the
relative free assertion, including continuity of the attaching maps.
\end{proof}


\begin{theorem}[Continuous comparison]\label{thm:comparison}
Let $X_\bullet\to G$ be an augmented based free simplicial object
in $\C$, and let $Y_\bullet\to H$ be a resolution. Every morphism
$\alpha:G\to H$ extends to a simplicial morphism
$f:X_\bullet\to Y_\bullet$. Any two extensions $f_0,f_1$ are
simplicially homotopic over $\alpha$, with all homotopy components
morphisms in $\C$.
Consequently any two based free resolutions of the same group are
simplicially homotopy equivalent over that group. No finite-rank or
finite-presentation assumption is needed.
\end{theorem}

\begin{proof}
Apply Lemma~\ref{lem:extension} to the relative free extension
$1\to X$ to obtain $f$. For two given extensions, their coproduct
defines $X\amalg X\to Y$ over $\alpha$. By
Lemma~\ref{lem:cylinder} and Lemma~\ref{lem:extension} it extends to
\[
 H:C(X)\longrightarrow Y,\qquad H i_0=f_0,\quad H i_1=f_1.
\]
Here is the explicit passage to the usual simplicial homotopy.
Write $H_{n,\tau}:X_n\to Y_n$ for the restriction to the
$\tau$-summand. For $0\leq i\leq n$ let
$\tau_i:[n+1]\to[1]$ be $0$ on $0,\ldots,i$ and $1$ on
$i+1,\ldots,n+1$, and put
\[
 h_i^n=H_{n+1,\tau_i}s_i:X_n\longrightarrow Y_{n+1}.
\]
All these maps are morphisms in $\C$. The identities follow directly
by precomposing the labels $\tau_i$ and using the simplicial identities
in $X$. More explicitly, they are
\begin{align*}
 d_0h_0&=f_1,&d_{n+1}h_n&=f_0,\\
 d_jh_i&=h_{i-1}d_j &&(j<i),\\
 d_jh_i&=h_i d_{j-1} &&(j>i+1),\\
 d_{i+1}h_i&=d_{i+1}h_{i+1} &&(i<n),\\
 s_jh_i&=h_{i+1}s_j &&(j\leq i),\\
 s_jh_i&=h_i s_{j-1} &&(j>i).
\end{align*}
These give a homotopy from $f_0$ to $f_1$ with the indicated endpoint
convention. The augmentation is preserved because $H$ is over
$\alpha$.

For two based free resolutions of $G$, construct maps in both
directions over $\id_G$. Each composite and the relevant identity
are extensions of $\id_G$. The homotopy uniqueness just proved
therefore makes the maps homotopy inverses.
\end{proof}

\begin{thebibliography}{99}
\bibitem[Inassaridze]{Inassaridze} H.~Inassaridze, \emph{Non-Abelian Homological Algebra and Its Applications}, Mathematics and Its Applications \textbf{421}, Kluwer Academic Publishers, Dordrecht, 1997. Chapter~2, ``Derived functors'', pp.~35--78. \href{https://doi.org/10.1007/978-94-015-8853-9}{doi:10.1007/978-94-015-8853-9}.
\bibitem[Keune]{Keune} F.~Keune, Derived functors and algebraic $K$-theory, in H.~Bass (ed.), \emph{Algebraic K-Theory I}, Lecture Notes in Mathematics \textbf{341}, Springer, 1973, 166--176.
\bibitem[RibesZalesskii]{RibesZalesskii} L.~Ribes and P.~Zalesskii, \emph{Profinite Groups}, second edition, Ergebnisse der Mathematik und ihrer Grenzgebiete (3) \textbf{40}, Springer, 2010. \href{https://doi.org/10.1007/978-3-642-01642-4}{doi:10.1007/978-3-642-01642-4}.
\end{thebibliography}
\end{document}
